% Clase 4 — Introducción a Búsqueda Informada
% Lectura: AIMA Sección 3.5
\section{Búsqueda Informada: Heurísticas}

\subsection{Función Heurística}

\begin{definition}[Heurística]
$h(n)$ es una función que estima el costo desde el nodo $n$ hasta la meta más cercana. Debe ser \textbf{admisible}: nunca sobreestima el costo real.
\end{definition}

\begin{remark}
Una heurística es \textbf{admisible} si $h(n) \leq h^*(n)$ para todo $n$, donde $h^*(n)$ es el costo real al llegar a la meta.
\end{remark}

\subsection{Greedy Best-First Search}

Expande el nodo que parece más cercano a la meta según $h(n)$.
\begin{itemize}
    \item Función de evaluación: $f(n) = h(n)$
    \item \textbf{No es óptimo} ni completo en general.
    \item Rápido en la práctica, puede quedar atrapado en mínimos locales.
\end{itemize}

\subsection{A*}

Combina el costo real del camino y la heurística:
$$f(n) = g(n) + h(n)$$

\begin{itemize}
    \item $g(n)$: costo del camino desde inicio hasta $n$
    \item $h(n)$: estimación del costo desde $n$ hasta la meta
    \item \textbf{Completo} y \textbf{óptimo} si $h$ es admisible (búsqueda en árbol) o consistente (búsqueda en grafo).
\end{itemize}

\begin{definition}[Consistencia / Monotonía]
$h$ es consistente si para todo nodo $n$ y sucesor $n'$:
$$h(n) \leq c(n, a, n') + h(n')$$
La consistencia implica admisibilidad.
\end{definition}

\begin{example}[Heurísticas comunes en una cuadrícula]
    Para un agente que se mueve en una cuadrícula 2D desde $(x_1,y_1)$ hasta $(x_2,y_2)$:
    \begin{itemize}
        \item \textbf{Distancia Manhattan}: $h(n) = |x_1-x_2| + |y_1-y_2|$. Admisible si el agente solo se mueve en
        horizontal/vertical (no sobreestima, porque es exactamente el número mínimo de pasos sin obstáculos).
        \item \textbf{Distancia euclidiana}: $h(n) = \sqrt{(x_1-x_2)^2 + (y_1-y_2)^2}$. Admisible si se permite
        movimiento diagonal, ya que nunca es mayor que el camino real.
    \end{itemize}
    Usar Manhattan cuando solo hay 4 movimientos posibles evita sobreestimar; usarla con movimiento diagonal
    permitido violaría la admisibilidad, porque el camino real podría ser más corto que la suma de catetos.
\end{example}

\begin{figure}[H]
\centering
\begin{tikzpicture}[
  nodo/.style={draw=black, thick, circle, minimum size=11mm, font=\scriptsize},
  descartado/.style={draw=black, thick, dashed, circle, minimum size=11mm, font=\scriptsize},
  meta/.style={draw=black, thick, double, circle, minimum size=11mm, font=\scriptsize}
]
\node[nodo] (arad)     at (0,0)    {\shortstack{A\\$g0,h366$}};
\node[nodo] (sibiu)    at (3,0.9)  {\shortstack{S\\$g140,h253$}};
\node[nodo] (zerind)   at (3,-0.9) {\shortstack{Z\\$g75,h374$}};
\node[nodo] (fagaras)  at (6,2.0)  {\shortstack{F\\$g239,h178$}};
\node[nodo] (rimnicu)  at (6,0.2)  {\shortstack{R\\$g220,h193$}};
\node[descartado] (bfagaras) at (9,2.0)  {\shortstack{B\\$g450,h0$}};
\node[nodo] (pitesti) at (9,0.2)  {\shortstack{P\\$g317,h100$}};
\node[meta] (bpitesti) at (12,0.2) {\shortstack{B\\$g418,h0$}};

\draw[-{Latex},thick] (arad) -- (sibiu);
\draw[-{Latex},thick] (arad) -- (zerind);
\draw[-{Latex},thick] (sibiu) -- (fagaras);
\draw[-{Latex},thick] (sibiu) -- (rimnicu);
\draw[-{Latex},thick,dashed] (fagaras) -- (bfagaras);
\draw[-{Latex},thick] (rimnicu) -- (pitesti);
\draw[-{Latex},thick] (pitesti) -- (bpitesti);
\end{tikzpicture}
\caption{Traza de A* sobre el mapa de Rumania (heurística: distancia en línea recta a Bucarest). A* siempre expande
el nodo de menor $f(n)=g(n)+h(n)$ de la frontera. Aunque Fagaras (F) llega primero a Bucarest con $f=450$ (nodo
punteado: se genera pero \textbf{se descarta}), Rimnicu (R) tiene $f=220+193=413$, menor que $f(F)=239+178=417$, así
que A* explora primero $R \to$ Pitesti (P) con $f=317+100=417$. Al final, la meta alcanzada por Pitesti tiene
$f=418$, menor que los $450$ de la rama de Fagaras, así que esa es la \textbf{solución óptima} (nodo de doble borde).}
\end{figure}
